\medskip

L'objectif de cet exercice est d'étudier la suite $\left(u_{n}\right)$ définie pour tout entier naturel $n$ par :

\[\left\{\begin{array}{l}
u_{0}=0 \\
u_{1}=\dfrac{1}{2} \\
u_{n+2}=u_{n+1}-\dfrac{1}{4} u_{n}
\end{array}\right.\]

\medskip

\textbf{Partie A : Conjecture}

\medskip

\begin{enumerate}
\item Voici le tableau complété (on peut calculer rapidement les termes à la main, puis vérifier à la calculatrice) :

\begin{center}
\begin{tabularx}{0.8\linewidth}{*{7}{|>{\centering\arraybackslash}X}|}\hline
		$n$ & 0 & 1 & 2 & 3 & 4 & 5 \\ \hline
		$u_{n}$ \rule{0pt}{18pt}& 0 & $\dfrac{1}{2}$ & $\dfrac{1}{2}$ & $\dfrac{3}{8}$ & $\dfrac{1}{4}$ & $\dfrac{5}{32}$ \\ [7pt]\hline
\end{tabularx}
\end{center}

\item Par exploration à la calculatrice, les termes de la suite semblent décroitrent, tout en restant strictement positifs, on a $u_{100}\approx8\times 10^{-29}$, et $u_{1000}\approx9\times 10^{-299}$.

On suppose que la suite converge vers 0.
\end{enumerate}

\medskip

\textbf{Partie B : Étude d'une suite auxiliaire}

\medskip

\begin{enumerate}
\item On a : \quad$w_{0} = u_1 - \dfrac{1}{2}u_0 = \dfrac{1}{2} - \dfrac{1}{2}\times 0 = \dfrac{1}{2}$.
\item On va établir la relation de récurrence de $(w_n)$. Soit $n$ un entier naturel.

$\begin{aligned}
	w_{n+1} &= u_{(n+1) + 1} - \dfrac{1}{2} u_{(n+1)} \quad \text{en appliquant la définition de $w$ au rang }(n+1)\\
	&= u_{n+2} - \dfrac{1}{2}u_{n+1}\\
	&=\left(u_{n+1} - \dfrac{1}{4}u_n\right)	- \dfrac{1}{2}u_{n+1}\quad\text{en appliquant la relation de récurrence de }u. \\
	&=\dfrac{1}{2}u_{n+1} - \dfrac{1}{4}u_n\\
	&=\dfrac{1}{2}\left(u_{n+1} - \dfrac{1}{2}u_n\right) \\
	&=\dfrac{1}{2}w_n\quad \text{en appliquant la définition de $w$ au rang }n\\
\end{aligned}$

Ainsi, on a : \quad $\forall n \in \N,\quad w_{n+1} = \dfrac{1}{2}w_n$.

Cette relation de récurrence établit que $(w_n)$ est une suite géométrique, de raison $q = \dfrac{1}{2}$, et de premier terme $w_0 = \dfrac{1}{2}$.

\item Puisque la suite est géométrique, on a la propriété classique :

$\forall n \in \N,\quad w_n = w_0 \times q^n = \dfrac{1}{2} \times \left(\dfrac{1}{2}\right)^n = \left(\dfrac{1}{2}\right)^{n+1}$.
\item Soit $n$ un entier naturel. On reprend la définition de $(w_n)$ :

$\begin{aligned}
	w_n = u_{n+1} - \dfrac{1}{2}u_n&\iff
	\left(\dfrac{1}{2}\right)^{n+1} = u_{n+1} - \dfrac{1}{2}u_n\quad \text{d'après l'expression explicite de }w_n\\
	&\iff u_{n+1}=\left(\dfrac{1}{2}\right)^{n+1}+\dfrac{1}{2} u_{n}
\end{aligned}$

On arrive bien à la relation de récurrence demandée.

\item Pour tout $n$ entier naturel, on pose : $P_n$ l'affirmation : \og $u_n = n\left(\dfrac{1}{2}\right)^{n}$\fg{}.

\textbf{Initialisation :} On a d'une part $u_0 = 0$ et, d'autre part : $0\times\left(\dfrac{1}{2}\right)^0 = 0\times 1 = 0 $.

L'affirmation est donc vraie au rang 0.

\textbf{Hérédité :} Pour un entier naturel $n$ donné, on suppose que la propriété $P_n$ est vraie, c'est-à-dire : $u_n = n\left(\dfrac{1}{2}\right)^n$.

On a : \quad
$\begin{aligned}[t]
u_{n+1} &= \left(\dfrac{1}{2}\right)^{n+1} + \dfrac{1}{2}u_n \quad \text{d'après la relation de récurrence de la question \textbf{B. 4}}\\
	& = \left(\dfrac{1}{2}\right)^{n+1} + \dfrac{1}{2}\times n\left(\dfrac{1}{2}\right)^{n}\quad\text{par hypothèse de récurrence}\\
	& = \left(\dfrac{1}{2}\right)^{n+1} + n\left(\dfrac{1}{2}\right)^{n+1}\\
	& = \left(\dfrac{1}{2}\right)^{n+1}\times\big(1 +n\big)\\
	u_{n+1}&=(n+1)\left(\dfrac{1}{2}\right)^{n+1} \quad \text{c'est l'affirmation }P_{n+1}
\end{aligned}$

\textbf{Conclusion :} L'affirmation $P_0$ est vraie, et, pour tout entier naturel $n$, la véracité de l'affirmation $P_n$ est héréditaire, donc, par principe de récurrence :

$\forall n \in \N,\quad u_n = n\left(\dfrac{1}{2}\right) ^{n}$.
\end{enumerate}

\medskip

\textbf{Partie C : Étude de la suite \boldmath$(u_{n})$\unboldmath}

\medskip

\begin{enumerate}
\item Soit $n$ un entier naturel non nul, donc supérieur ou égal à 1 :

$\begin{aligned}
			u_{n+1} - u_n &= (n+1)\left(\dfrac{1}{2}\right)^{n+1} - n\left(\dfrac{1}{2}\right)^{n}\quad \text{d'après la question \textbf{B. 5.}}\\
	&= (n+1)\left(\dfrac{1}{2}\right)^{n+1} - n\left(\dfrac{1}{2}\right)^{n+1} \times 2\\
	&= \left(\dfrac{1}{2}\right)^{n+1} \times\big((n+1) -2n \big) \\
	&= \left(\dfrac{1}{2}\right)^{n+1} \times\big(1-n\big) \\
\end{aligned}$

La différence $u_{n+1} - u_n$ est égale au produit de deux nombres de signe contraire, car, pour $n$ entier naturel supérieur ou égal à 1 :
\begin{itemize}
			\item $\left(\dfrac{1}{2}\right)^{n+1}$ est positif strictement;
			\item $(1 - n)$ est négatif ou nul
\end{itemize}

La différence $u_{n+1} - u_n$ est donc négative ou nulle pour tout $n$ supérieur ou égal à 1, on en déduit donc que la suite $(u_{n})$ est décroissante à partir du rang $n=1$.

\item L'expression du terme général de la suite $(u_n)$ permet d'affirmer que la suite est minorée par 0, car chaque terme est le produit de $n$, entier naturel, donc positif et de $\left(\dfrac{1}{2}\right)^{n}$, strictement positif, car $\dfrac{1}{2}$ est strictement positif.

De plus, la suite est décroissante, à partir du rang $n= 1$.

La suite est donc décroissante (à partir du rang $n=1$) et minorée par 0 : on en déduit qu'elle converge, vers une limite $\ell$ dont on sait que $\ell \geqslant0$.

\item On admet que la limite de la suite $\left(u_{n}\right)$ est solution de l'équation : $\ell=\ell-\dfrac{1}{4} \ell$.

Résolvons cette équation :\quad 
$\begin{aligned}[t]
	\ell=\ell-\dfrac{1}{4} \ell
	&\iff \ell = \dfrac{3}{4}\ell\\
	&\iff \ell - \dfrac{3}{4}\ell = 0\\
	&\iff \dfrac{1}{4}\ell = 0\\
	&\iff \ell = 0\quad \text{car } \dfrac{1}{4}\neq 0\\
\end{aligned}$

L'équation ayant une unique solution, puisque la limite doit être une solution de l'équation, on a donc la limite de la suite $\left(u_{n}\right)$ qui est 0, l'unique solution de l'équation.

Cela vient confirmer notre conjecture de la \textbf{partie A}.
	\end{enumerate}
